\documentclass[12pt,a4paper]{article}

\usepackage[utf8]{inputenc}
\usepackage[T2A]{fontenc}
\usepackage[russian]{babel}
\usepackage{amsmath,amssymb}
\usepackage{geometry}
\usepackage{graphicx}
\usepackage{float}

\geometry{margin=2.2cm}

\newcommand{\plotplaceholder}[2][0.9\textwidth]{%
    \IfFileExists{#2}{%
        \includegraphics[width=#1]{#2}%
    }{%
        \fbox{%
            \begin{minipage}[c][0.22\textheight][c]{#1}
            \centering
            Заглушка для графика\\[4pt]
            \texttt{\detokenize{#2}}
            \end{minipage}%
        }%
    }%
}

\begin{document}

\title{Метод Монте-Карло}
\author{}
\date{}
\maketitle

\section*{Цель}

В первом разделе была аналитически получена нормированная ОИСКО ядерной
оценки с прямоугольным ядром. Метод Монте-Карло позволяет проверить эту
формулу численно: мы многократно генерируем выборки, для каждой выборки
строим оценку плотности и считаем её интегральную квадратическую ошибку.
После усреднения по большому числу повторов результат должен быть близок к
аналитической RMISE.

Истинная плотность остаётся прежней:
\[
a=\sqrt{3},\qquad
f(x)=\frac{1}{2a}\mathbf 1_{\{|x|\le a\}}.
\]
Ядерная оценка плотности имеет вид
\[
\widehat f_{\xi}(x)
=
\frac{1}{n\xi}\sum_{j=1}^n
K\left(\frac{x-X_j}{\xi}\right),
\]
где
\[
K(x)=\frac{1}{2a}\mathbf 1_{\{|x|\le a\}}.
\]

\section*{Алгоритм Монте-Карло}

Фиксируем объём выборки $n$ и параметр сглаживания $\xi$. Один прогон
Монте-Карло состоит из следующих шагов.

\begin{enumerate}
\item Сгенерировать выборку
\[
X_1,\ldots,X_n\sim U[-a,a].
\]
\item Построить ядерную оценку $\widehat f_{\xi,s}(x)$ по этой выборке.
\item Посчитать интегральную квадратическую ошибку
\[
ISE_s(\xi)
=
\int_{-\infty}^{\infty}
\bigl(\widehat f_{\xi,s}(x)-f(x)\bigr)^2\,dx.
\]
\item Нормировать её:
\[
ISE_s^*(\xi)
=
\frac{ISE_s(\xi)}{\|f\|_2^2}
=
2a\,ISE_s(\xi).
\]
\end{enumerate}

После $N_{MC}$ независимых повторов получаем оценку средней ошибки:
\[
\widehat{\operatorname{RMISE}}_{MC}(\xi)
=
\frac{1}{N_{MC}}\sum_{s=1}^{N_{MC}} ISE_s^*(\xi).
\]
Именно эта величина сравнивается с аналитической формулой
$\operatorname{RMISE}(\xi)$.

\section*{Численная реализация}

Интеграл в $ISE_s(\xi)$ считается численно на равномерной сетке по $x$.
Так как носитель ядерной оценки шире носителя истинной плотности, сетка
берётся на отрезке, покрывающем все возможные значения, где оценка может
быть ненулевой. Интегрирование выполняется методом трапеций.

В расчётах использованы параметры
\[
N_{MC}=3000,\qquad
n\in\{30,100\}.
\]
Значения $\xi$ взяты как около оптимума, так и в стороне от него, чтобы
проверить всю форму кривой RMISE.

\section*{Результаты}

На графиках ниже синяя линия --- аналитическая формула, а красные точки ---
оценка методом Монте-Карло. Вертикальные отрезки показывают приближённый
95\%-й доверительный интервал для среднего значения Монте-Карло.

\begin{figure}[H]
\centering
\plotplaceholder[0.82\textwidth]{mc_kernel_rmise_n30.png}
\caption{Сравнение аналитической RMISE и Монте-Карло при $n=30$.}
\end{figure}

\begin{figure}[H]
\centering
\plotplaceholder[0.82\textwidth]{mc_kernel_rmise_n100.png}
\caption{Сравнение аналитической RMISE и Монте-Карло при $n=100$.}
\end{figure}

Для количественной проверки рассмотрим относительное отклонение
\[
\frac{
\widehat{\operatorname{RMISE}}_{MC}(\xi)
-
\operatorname{RMISE}(\xi)
}
{\operatorname{RMISE}(\xi)}.
\]

\begin{figure}[H]
\centering
\plotplaceholder[0.82\textwidth]{mc_relative_error.png}
\caption{Относительное отклонение оценки Монте-Карло от аналитической формулы.}
\end{figure}

Отклонения имеют случайный характер и остаются малыми по сравнению с самим
значением RMISE. Это подтверждает аналитическую формулу, полученную ранее
для прямоугольного ядра.

\end{document}
